\scanpage{123}
% archive foliation: 122
OPS \deleted{\unsure{quitte}} par changement $\unsure{\alpha\in}\deleted{\unsure{A}}A^p$, puis $\alpha\unsure{=|\ne}0\deleted{\unsure{\alpha}}$. \deleted{\unsure{alors}} $t\in\unsure{\mathfrak r B|\sqrt{0_B}}$, donc
\[
D^{(1)}s=a_1+a_2t+\cdots\equiv a_1\ (\unsure{\mathfrak r B}),
\]
donc $D^{(1)}s\in B^*\Longleftrightarrow a_1\in A^*$. OPS, quitte à remplacer $s$ par $a_1^{-1}(s-a_0)$, que $s=t+a_2t^2+\cdots+a_{p-1}t^{p-1}$ puis, quitte à faire encore un chgt de variable \deleted{\unsure{et}}
\[
s'=\frac{s}{1+a_2s}=\deleted{\unsure{\frac1{1+a_2s^{-1}}}}=\deleted{\unsure{s}}\sum_{n\ge0}\inserted{(-1)^n}a_2^n s^{n+1}
\]
$=s\unsure{-|+}a_2s^2+\cdots\equiv t\bmod t^3$ donc $s$ de la forme
\[
s=t+\text{termes d'ordre }\ge3.
\]
\deleted{Alors} À prouver que, dans ces conditions, la condition
\[
c_0(s)=0\Longrightarrow s=t.
\]
\deleted{\unsure{Si on a}} Il suffit de prouver que si
\[
s=t+at^d+\cdots\quad(\text{termes d'ordre }\ge d+1)
\]
où $\deleted{d\ge3}\inserted{3\le d\le p-1}$, alors $a=0$.
\[
D^{(3)}s\equiv a\binom d3t^{d-3}\bmod(t^{d-2})
\]
\[
D^{(2)}s\equiv a\binom d2t^{d-2}\bmod(t^{d-1})\quad\text{et a fortiori mod }t^{d-2}
\]
\[
D^{(1)}s\equiv1\deleted{+ad\,t^{d-1}}\bmod(t^{d-1})\quad\text{et a fort. mod }t^{d-2}
\]
\[
D^{(3)}s\,D^{(1)}s\equiv a\binom d3t^{d-3}\deleted{+a^2d\binom d3t^{2d-4}}\quad\bmod t^{d-2}
\]
\[
(D^{(2)}s)^2\equiv0\quad\bmod t^{d-2}
\]
donc $a\binom d3t^{d-3}\equiv0\quad\bmod t^{d-2}$

donc $a\binom d3=0$ dans $A$,

i.e. $\dfrac{d(d-1)(d-2)}6\,a=0$ dans $A$

or 6 inv. dans $A$, et aussi $d(d-1)(d-2)$ car $3\le d\le p-1$ et $A$ de car. $p>3$, donc $a=0$ OK.
